\chapter{Compactificaciones y espacios universales}\label{cap:nuevo07}\label{chap:16}

\section{Añadir límites conservando la topología}\label{lectura:16}
Compactificar significa situar un espacio dentro de otro compacto, añadiendo puntos que pueden representar comportamientos límite. La inclusión debe conservar exactamente la topología original: una inyección continua por sí sola puede cambiarla. Desarrollamos la lectura 16 de \cite{valencia}, corrigiendo esa distinción y resolviendo sus ejercicios, incluso cuando el enunciado original requiere modificación. Necesitaremos productos numerables, compacidad y separabilidad, junto con la conexidad y los caminos estudiados en el capítulo \ref{cap:nuevo06}. El ejercicio \ref{ej:lectura-16-2} incluye una introducción autónoma al índice de un lazo para justificar las diferencias topológicas que solicita.

\section{La definición y dos modelos compactos}
\begin{definicion}
Una \emph{compactificación métrica}\index{compactificacion@compactificación} de $X$ es un par $(K,j)$, donde $K$ es compacto métrico y $j:X\to K$ es un \emph{encaje}: un homeomorfismo de $X$ sobre $j(X)$, con $j(X)$ denso en $K$.
\end{definicion}
\begin{observacion}[Por qué exigir un encaje]
La definición original solo exige que $j$ sea continua e inyectiva. Eso no permite deducir que $X$ sea separable. Por ejemplo, tomemos $\R$ con la métrica discreta y una biyección $j:\R\to(0,1)$. Toda función desde un espacio discreto es continua, de modo que $j$ es una inyección continua de imagen densa en $[0,1]$, aunque el dominio discreto no numerable no es separable. Su inversa no es continua. Con la definición de encaje, el argumento correcto sí funciona: un compacto métrico es segundo numerable, sus subespacios son segundo numerables y todo espacio segundo numerable es separable, eligiendo un punto en cada miembro no vacío de una base numerable. Por tanto todo espacio con compactificación métrica debe ser separable.
\end{observacion}

\begin{definicion}
El \emph{conjunto de Cantor}\index{conjunto de Cantor} es
\[
C=\left\{\sum_{n=1}^\infty\frac{2b_n}{3^n}:b_n\in\{0,1\}\right\}\subset[0,1].
\]
El \emph{cubo de Hilbert}\index{cubo de Hilbert} es $H=[0,1]^{\N}$ con la métrica
\[
D(x,y)=\sum_{n=1}^\infty2^{-n}\frac{|x_n-y_n|}{1+|x_n-y_n|}.
\]
\end{definicion}
El uso del nombre cubo de Hilbert para el compacto $K\subset\ell^2$ del ejercicio \ref{ej:lectura-14-5} es compatible: el ejercicio \ref{ej:lectura-16-4} construye un homeomorfismo entre ambos modelos.

\begin{proposicion}\label{ch16:cantor}
$C$ y $H$ son compactos. El mapa
\[
\theta:\{0,1\}^{\N}\longrightarrow C,\qquad
\theta(b)=\sum_{n=1}^\infty2b_n3^{-n},
\]
es un homeomorfismo. Por tanto el mapa $\theta(b)\mapsto b$ encaja $C$ en $H$ como el conjunto de sus vértices.
\end{proposicion}
\begin{proof}
Los espacios $\{0,1\}^{\N}$ y $H$ son compactos por el ejercicio \ref{ej:lectura-14-1}. Si dos sucesiones coinciden hasta la coordenada $N$, sus imágenes por $\theta$ difieren a lo sumo en $\sum_{n>N}2/3^n=3^{-N}$, lo que prueba continuidad en la topología producto. El mapa es sobreyectivo por la definición de $C$. Si el primer dígito distinto está en la posición $k$, la diferencia de los términos de esa posición es $2/3^k$, mientras que toda la cola puede compensar a lo sumo $1/3^k$; las imágenes difieren por lo menos en $1/3^k$. Así $\theta$ es inyectivo. Una biyección continua de un compacto a un espacio métrico es un homeomorfismo, por el ejercicio \ref{ej:lectura-13-5}. En particular, su imagen $C$ es compacta.

Para relacionar la fórmula con la construcción por tercios, un prefijo $b_1,\ldots,b_N$ determina el intervalo
\[
\left[\sum_{n=1}^N\frac{2b_n}{3^n},\ 
\sum_{n=1}^N\frac{2b_n}{3^n}+3^{-N}\right].
\]
La unión de estos $2^N$ intervalos se obtiene eliminando los tercios centrales durante $N$ etapas. Los intervalos de cada etapa son disjuntos, y cada punto de todas las etapas determina prefijos compatibles; su distancia a la suma correspondiente es a lo sumo $3^{-N}$ para todo $N$. Por tanto la intersección de las etapas coincide exactamente con $C$.
\end{proof}

\section{El encaje universal del espacio separable}
\begin{teorema}\label{ch16:encaje}
Todo espacio métrico separable se encaja en $H$. En consecuencia, un espacio métrico admite una compactificación métrica si y solo si es separable.
\end{teorema}
\begin{proof}
Si $X=\varnothing$, la aplicación vacía $X\to H$ es un encaje, y la clausura vacía de su imagen es una compactificación de $X$. Para $X\ne\varnothing$, sustituimos $d$ por $\rho=d/(1+d)$, que es una métrica equivalente, acotada por $1$. Su desigualdad triangular se demostró en el ejercicio \ref{ej:lectura-14-1}; la equivalencia se sigue de que $t/(1+t)$ es continua, estrictamente creciente y tiene inversa $s/(1-s)$ para $0\le s<1$.

Sea $(a_n)$ una sucesión cuyo recorrido es denso; si el conjunto denso es finito, repetimos sus puntos. Definimos
\[
J(x)=(\rho(x,a_n))_{n\ge1}\in H.
\]
Cada coordenada es continua y, de hecho,
\[
D(J(x),J(y))\le\sum_{n\ge1}2^{-n}|\rho(x,a_n)-\rho(y,a_n)|\le\rho(x,y).
\]
Por tanto $J$ es continua. Si $J(x)=J(y)$, dado $\eps>0$ elegimos $a_n$ con $\rho(x,a_n)<\eps/2$. La igualdad de coordenadas implica $\rho(y,a_n)=\rho(x,a_n)$, así que $\rho(x,y)<\eps$. Como $\eps$ es arbitrario, $x=y$.

Falta la continuidad de la inversa, que no se sigue solamente de la inyectividad. Fijemos $x\in X$ y $\eps>0$. Elegimos $a_n$ con $\rho(x,a_n)<\eps/4$. Si $J(y)$ está en una vecindad de $J(x)$ que haga variar la coordenada $n$ menos que $\eps/4$, entonces $\rho(y,a_n)<\eps/2$ y
\[
\rho(y,x)\le\rho(y,a_n)+\rho(a_n,x)<3\eps/4<\eps.
\]
Tal vecindad existe porque la evaluación de una coordenada es continua en la métrica producto. Esto prueba que $J^{-1}$ es continua sobre $J(X)$. Finalmente, $K=\cl{J(X)}^{\,H}$ es compacto como cerrado de $H$, y $J(X)$ es denso en $K$. La necesidad de separabilidad se probó después de la definición.
\end{proof}

\begin{ejemplo}[Dos extremos para la recta]
La función
\[
u(x)=\frac12\left(1+\frac{x}{1+|x|}\right)
\]
es un homeomorfismo de $\R$ sobre $(0,1)$. Es continua, estrictamente creciente y su inversa se obtiene poniendo $s=2t-1$ y $x=s/(1-|s|)$. Como $(0,1)$ es denso en $[0,1]$, este último compactifica la recta. Si añadimos los símbolos $-\infty,+\infty$ y ponemos $u(-\infty)=0$, $u(+\infty)=1$, la fórmula $\rho(x,y)=|u(x)-u(y)|$ metriza la recta extendida y la identifica con $[0,1]$. La métrica extendida induce en la recta original su topología usual, aunque no su distancia usual.

Se pueden adoptar convenciones de cálculo con infinitos, como $x+(+\infty)=+\infty$ para $x\in\R$ y $x(+\infty)=+\infty$ para $x>0$. Estas convenciones no convierten la recta extendida en un cuerpo; $+\infty-\infty$ queda sin definir. La convención $0\cdot\infty=0$, útil en ciertos contextos de medida, tampoco es una extensión continua de la multiplicación.
\end{ejemplo}

\begin{ejemplo}[El espacio de Baire dentro de Cantor]
Sea $\mathcal N=\N^{\N}$, con $\N=\{1,2,\ldots\}$, y distancia $0$ entre sucesiones iguales y $1/k$ si el primer desacuerdo aparece en la posición $k$. Codificamos $x=(x_n)$ por la sucesión binaria
\[
\underbrace{0\cdots0}_{x_1-1}1\,
\underbrace{0\cdots0}_{x_2-1}1\,
\underbrace{0\cdots0}_{x_3-1}1\cdots.
\]
Su imagen es el conjunto de sucesiones binarias con infinitos unos. La codificación es biyectiva sobre esa imagen porque las longitudes de los bloques recuperan $x_n-1$. Es continua: fijar los primeros $m$ números fija un prefijo binario de longitud al menos $m$. La inversa es continua: para recuperar los primeros $m$ números basta leer el prefijo hasta el $m$-ésimo uno. La imagen es densa en $\{0,1\}^{\N}$, ya que cualquier prefijo finito puede prolongarse con infinitos unos. Mediante $\theta$, obtenemos una compactificación de $\mathcal N$ homeomorfa a $C$. El ajuste $x_n-1$ evita que la convención $\N\ni1$ prohíba unos consecutivos; si se usan naturales desde cero, se pueden tomar directamente $x_n$ ceros.\index{espacio de Baire}
\end{ejemplo}

\section{Todo compacto métrico no vacío es imagen de Cantor}
\begin{teorema}\label{ch16:imagenCantor}
Un espacio métrico no vacío $X$ es compacto si y solo si existe una aplicación continua sobreyectiva $C\to X$.
\end{teorema}
\begin{proof}
La suficiencia es el teorema de imagen continua. Para la necesidad construiremos una dirección de aproximación para cada sucesión de dígitos, de modo que sus aproximaciones tengan un único límite.

Partimos de $K_{\varnothing}=X$. Cubrimos cada compacto no vacío $K_s$ ya construido por finitos conjuntos compactos no vacíos contenidos en él y de diámetro a lo sumo $2^{-n}$ en el nivel $n$. Esto se consigue intersectando $K_s$ con finitas bolas cerradas de radio $2^{-n-1}$ cuyos centros están en $K_s$ y cuyas bolas abiertas cubren $K_s$. En cada nivel hay solo finitos nodos: sea $m_n$ el mayor número de hijos de un nodo del nivel anterior. Si a un nodo le faltan hijos para llegar a $m_n$, repetimos uno de los existentes. Así los hijos se etiquetan uniformemente por $1,\ldots,m_n$.

Para $a=(a_n)\in\prod_n\{1,\ldots,m_n\}$, la cadena
\[
K_{a_1}\supset K_{a_1a_2}\supset\cdots
\]
está formada por compactos no vacíos. Tiene intersección no vacía por la propiedad de intersección finita dentro de $X$, y la intersección tiene a lo sumo un punto porque los diámetros tienden a cero. Llamemos $P(a)$ a ese único punto. Si dos sucesiones coinciden hasta el nivel $n$, sus imágenes están en un mismo conjunto de diámetro a lo sumo $2^{-n}$; esto prueba continuidad de $P$. Para cualquier $x\in X$, elegimos sucesivamente un hijo que contenga $x$, lo cual es posible porque los hijos cubren al padre. Se obtiene $a$ con $P(a)=x$, de modo que $P$ es sobreyectiva.

Falta obtener este producto finito a partir de dígitos binarios. Elegimos enteros $l_n\ge1$ con $2^{l_n}\ge m_n$ y una sobreyección $q_n:\{0,1\}^{l_n}\to\{1,\ldots,m_n\}$. Dividimos una sucesión binaria en bloques sucesivos de longitudes $l_1,l_2,\ldots$ y aplicamos $q_n$ al bloque $n$. El mapa resultante $Q$ es continuo porque cada coordenada depende de finitos dígitos, y es sobreyectivo porque podemos elegir por separado una preimagen para cada coordenada. La composición $P\circ Q\circ\theta^{-1}:C\to X$ es la aplicación buscada.
\end{proof}

\begin{observacion}[La aplicación binaria y un salto lógico que hay que evitar]
La fórmula correcta $g(b)=\sum_{n\ge1}b_n/2^n$ lleva continuamente $\{0,1\}^{\N}$ sobre $[0,1]$. La continuidad sigue del control de la cola; la sobreyectividad se obtiene de la expansión binaria: para $t<1$, el algoritmo $b_n=\lfloor2t_{n-1}\rfloor$, $t_n=2t_{n-1}-b_n$, comenzando con $t_0=t$, da $t=\sum_{n\le N}b_n2^{-n}+2^{-N}t_N$, y el resto tiende a cero. Para $t=1$ se toman todos los dígitos iguales a uno. Con denominador $2^{n+1}$ y el índice inicial $n=1$, la imagen solo alcanzaría $[0,1/2]$.

Una biyección entre $\N$ y $\N\times\N$ reordena las coordenadas y prueba que $(\{0,1\}^{\N})^{\N}$ es homeomorfo a $\{0,1\}^{\N}$. La continuidad en ambas direcciones sigue de que cada coordenada se convierte en otra coordenada. Aplicando $g$ fila por fila obtenemos una sobreyección continua $C\to H$. Sin embargo, de $X\subset H$ no se deduce automáticamente una sobreyección $C\to X$: un mapa sobre $H$ no tiene por qué tomar valores en $X$. La construcción del teorema \ref{ch16:imagenCantor} cubre exactamente ese paso.
\end{observacion}

\section{Ejercicios resueltos: compactificaciones y obstrucciones topológicas}
\begin{ejercicio}[Lectura 16, ejercicio 1]\label{ej:lectura-16-1}
Probar que $x\mapsto e^{2i\arctan x}$ encaja $\R$ densamente en $S^1$. Probar que $S^1$ y $[0,1]$ no son homeomorfos.
\end{ejercicio}
\begin{solucion}
La función $2\arctan$ es un homeomorfismo de $\R$ sobre $(-\pi,\pi)$, y $t\mapsto e^{it}$ es un homeomorfismo de ese intervalo sobre $S^1\setminus\{-1\}$. Su inversa es el argumento en $(-\pi,\pi)$; es continuo en la circunferencia cortada porque localmente puede elegirse un intervalo angular de longitud menor que $2\pi$. Así la composición es un encaje. Su imagen es densa: el único punto que falta, $-1$, es límite de $e^{2i\arctan n}$.

Al quitar cualquier punto de $S^1$ queda un espacio homeomorfo a $\R$, por una rotación seguida de la parametrización anterior; por tanto queda conexo. Al quitar $1/2$ de $[0,1]$, quedan los dos abiertos relativos no vacíos $[0,1/2)$ y $(1/2,1]$, que son disjuntos. Un homeomorfismo induciría otro entre los espacios con los puntos correspondientes eliminados, y preservaría conexidad. Esto es imposible. La conexidad de los intervalos quedó demostrada en el lema \ref{ch17:intervalos} a partir de la propiedad del supremo.
\end{solucion}

\begin{ejercicio}[Lectura 16, ejercicio 2]\label{ej:lectura-16-2}
Probar que $S^2$, el toro $S^1\times S^1$ y $[0,1]^2$ son compactificaciones de $\R^2$ y son dos a dos no homeomorfos.
\end{ejercicio}
\begin{solucion}
\emph{Construcción de las compactificaciones.} Los tres espacios son compactos: la esfera es cerrada y acotada en $\R^3$, y los otros dos son productos finitos de compactos. La proyección estereográfica inversa
\[
\sigma(u,v)=\left(\frac{2u}{1+u^2+v^2},
\frac{2v}{1+u^2+v^2},\frac{u^2+v^2-1}{1+u^2+v^2}\right)
\]
es un homeomorfismo de $\R^2$ sobre $S^2\setminus\{(0,0,1)\}$. Sustituir la fórmula verifica que la suma de cuadrados es uno, y la inversa es $(x,y,z)\mapsto(x/(1-z),y/(1-z))$. Ambos mapas son continuos en sus dominios. El punto norte es límite de $\sigma(n,0)$, de modo que la imagen es densa. El producto de dos copias del encaje del ejercicio \ref{ej:lectura-16-1} encaja $\R^2$ densamente en el toro. El producto del mapa $u:\R\to(0,1)$ del ejemplo sobre la recta extendida encaja $\R^2$ densamente en $[0,1]^2$. La densidad en los productos se verifica en rectángulos abiertos, que forman una base.

El producto $S^1\times S^1$ está naturalmente en $\R^4$, no literalmente en $\R^3$. Su modelo usual en $\R^3$ es la imagen de
\[
(e^{i\alpha},e^{i\beta})\longmapsto
((2+\cos\beta)\cos\alpha,(2+\cos\beta)\sin\alpha,\sin\beta).
\]
Es continuo e inyectivo: el radio cilíndrico $\rho=\sqrt{x^2+y^2}$ determina $\cos\beta=\rho-2$, la tercera coordenada determina $\sin\beta$, y $(x/\rho,y/\rho)$ determina $e^{i\alpha}$. Como el dominio es compacto, es un homeomorfismo sobre su imagen.

\emph{Herramienta elemental para distinguirlos.} Un \emph{lazo} en un espacio es una aplicación continua $\gamma:[0,1]\to X$ con $\gamma(0)=\gamma(1)$. Decimos que se contrae si existe un mapa continuo $F:[0,1]^2\to X$ con $F(t,0)=\gamma(t)$, $F(t,1)$ constante en $t$ y $F(0,s)=F(1,s)$ para todo $s$. Esta propiedad es preservada por homeomorfismos. Demostraremos los hechos que necesitamos, sin suponer teoría del grupo fundamental. Usaremos la conexidad del intervalo, demostrada mediante el supremo en el lema \ref{ch17:intervalos}, y las nociones de caminos del capítulo \ref{cap:nuevo06}.\index{lazo}\index{homotopia@homotopía}

Todo camino $\gamma:[0,1]\to S^1$ admite un argumento continuo $\vartheta:[0,1]\to\R$ con $\gamma(t)=e^{i\vartheta(t)}$. Para construirlo, por continuidad uniforme elegimos una partición tal que, en cada subintervalo $[t_j,t_{j+1}]$, $|\gamma(t)-\gamma(t_j)|<1$. El cociente $\gamma(t)/\gamma(t_j)$ evita $-1$ y tiene un argumento principal continuo en $(-\pi,\pi)$. Elegida $\vartheta(0)$, definimos en ese subintervalo
\[
\vartheta(t)=\vartheta(t_j)+\operatorname{Arg}(\gamma(t)/\gamma(t_j)).
\]
Las definiciones se pegan en los extremos. Dos argumentos continuos difieren en un múltiplo constante de $2\pi$, pues su diferencia es continua y toma valores en el conjunto discreto $2\pi\Z$.

Para un lazo, el entero
\[
\operatorname{ind}(\gamma)=\frac{\vartheta(1)-\vartheta(0)}{2\pi}
\]
no depende del argumento elegido. Si dos lazos $\alpha,\beta$ satisfacen $\sup_t|\alpha(t)-\beta(t)|<2$, entonces $q(t)=\operatorname{Arg}(\beta(t)/\alpha(t))$ es continuo y $q(0)=q(1)$. Un argumento de $\beta$ se obtiene sumando $q$ a uno de $\alpha$, así que tienen el mismo índice. En una deformación continua $F(t,s)$, la continuidad uniforme permite dividir el intervalo de $s$ en finitos pasos tales que lazos sucesivos estén uniformemente a distancia menor que $1$. Por lo anterior el índice es constante durante la deformación. El lazo $e^{2\pi it}$ tiene índice $1$; un lazo constante tiene índice $0$. En consecuencia, el primero no se contrae en $S^1$.\index{indice de un lazo@índice de un lazo}

\emph{El toro tiene un lazo que no se contrae.} Consideremos $\gamma(t)=(e^{2\pi it},1)$. Si se contrajera en el toro, al componer la contracción con la primera proyección obtendríamos una contracción de $e^{2\pi it}$ en $S^1$, que acabamos de excluir.

\emph{Todos los lazos de la esfera se contraen.} Este paso necesita justificación: una curva continua puede llenar una región, así que no debemos suponer sin prueba que omite algún punto de la esfera. Sea $\gamma$ un lazo en $S^2$. Elegimos una partición tan fina que en cada intervalo $[t_j,t_{j+1}]$ se cumpla $\norm{\gamma(t)-\gamma(t_j)}<1/4$. Unimos los valores consecutivos $a_j=\gamma(t_j)$ y $a_{j+1}$ mediante el arco
\[
\beta(t)=\frac{(1-r)a_j+ra_{j+1}}{\norm{(1-r)a_j+ra_{j+1}}},\qquad
r=\frac{t-t_j}{t_{j+1}-t_j}.
\]
El denominador es positivo porque los extremos están a distancia menor que $1/4$, de modo que no son antipodales. Si $v=(1-r)a_j+ra_{j+1}$, entonces $\norm{v-a_j}<1/4$ y
\[
\norm{v/\norm v-a_j}\le |1-\norm v|+\norm{v-a_j}<1/2.
\]
Por tanto $\norm{\beta(t)-\gamma(t)}<3/4$. Los dos valores nunca son antipodales, y
\[
F(t,s)=\frac{(1-s)\gamma(t)+s\beta(t)}{\norm{(1-s)\gamma(t)+s\beta(t)}}
\]
define una deformación continua de lazos entre $\gamma$ y $\beta$.

La imagen de $\beta$ está en una unión finita de planos propios que pasan por el origen, intersectados con $S^2$, porque cada arco está en el plano generado por sus extremos; si ambos son iguales elegimos cualquier plano que contenga ese punto. Hay un punto $q\in S^2$ fuera de todos esos planos. Para verlo sin un argumento de medida, escribamos sus ecuaciones como $A_jx+B_jy+C_jz=0$, con algún coeficiente no cero. El vector $(1,t,t^2)$ pertenece al plano $j$ solo si $A_j+B_jt+C_jt^2=0$, ecuación con finitas raíces. Elegimos $t$ fuera de la unión finita de raíces y normalizamos el vector. Entonces $q$ no pertenece a la imagen de $\beta$.

Una rotación seguida de la proyección estereográfica identifica $S^2\setminus\{q\}$ con $\R^2$. Allí todo lazo se contrae por interpolación lineal con un punto fijo. Transportamos esa contracción a la esfera y la concatenamos con $F$, usando la primera mitad del parámetro para $F$ y la segunda para la contracción. Así $\gamma$ se contrae. La esfera y el toro no son homeomorfos.

\emph{El cuadrado y el toro tampoco son homeomorfos.} Todo lazo del cuadrado se contrae mediante $F(t,s)=(1-s)\gamma(t)+s(1/2,1/2)$, pues el cuadrado es convexo; el toro tiene el lazo no contraíble ya construido.

\emph{La esfera y el cuadrado no son homeomorfos.} Al quitar cualquier punto de la esfera queda un espacio homeomorfo a $\R^2$, donde todos los lazos se contraen. En cambio, al quitar $a=(1/2,1/2)$ del cuadrado, el lazo
\[
\gamma(t)=a+\tfrac14(\cos2\pi t,\sin2\pi t)
\]
no se contrae: componer una supuesta contracción con $x\mapsto(x-a)/\norm{x-a}$ daría una contracción del lazo de índice uno en $S^1$. Si existiera un homeomorfismo de la esfera al cuadrado, al eliminar la preimagen de $a$ produciría un homeomorfismo entre estos dos espacios perforados. La contradicción termina las tres comparaciones.
\end{solucion}

\begin{ejercicio}[Lectura 16, ejercicio 3: examen y corrección del enunciado]\label{ej:lectura-16-3}
El enunciado propone que el conjunto de Cantor sea una compactificación de todo espacio discreto infinito separable. Decidir si es válido con la definición de encaje y construir una compactificación correcta.
\end{ejercicio}
\begin{solucion}
Un espacio discreto separable es numerable, porque todo subconjunto denso debe contener cada singleton abierto. Sin embargo, \emph{no puede encajarse densamente en Cantor}. En efecto, si $D$ es un subespacio denso y discreto de un espacio métrico $K$, para cada $d\in D$ existe un abierto $U\subset K$ con $U\cap D=\{d\}$. Si $U$ contuviera un punto $y\ne d$, el abierto no vacío $U\setminus\{d\}$ tendría que tocar $D$ por densidad, contradicción. Así $U=\{d\}$ y $d$ es un punto aislado de $K$. Pero Cantor no tiene puntos aislados: manteniendo un prefijo binario y cambiando algún dígito posterior obtenemos puntos distintos arbitrariamente próximos.

La sugerencia de enumerar un subconjunto denso de vértices proporciona una inyección continua desde el discreto numerable, pero su inversa sobre la imagen no es continua. Esta es precisamente la dificultad de la definición original. Una versión correcta es la siguiente: si $X=\{x_1,x_2,\ldots\}$ es discreto, el mapa $x_n\mapsto1/n$ lo encaja densamente en
\[
K=\{0\}\cup\{1/n:n\ge1\}.
\]
Cada $1/n$ es aislado en $K$, de modo que la imagen tiene la topología discreta; $0$ es su único punto límite añadido. El conjunto $K$ es cerrado y acotado en $\R$, luego compacto.

La afirmación adicional de que todo espacio ultramétrico tiene a Cantor como compactificación tampoco es válida en esa generalidad: los discretos anteriores son ultramétricos, y los no separables ni siquiera admiten compactificación métrica. Sí puede afirmarse, por el teorema \ref{ch16:encaje}, que todo espacio ultramétrico \emph{separable} admite alguna compactificación métrica.
\end{solucion}

\begin{ejercicio}[Lectura 16, ejercicio 4]\label{ej:lectura-16-4}
Probar que $T:K\to H$, $T(x)_n=nx_n$, es un homeomorfismo, donde $K=\{x\in\ell^2:0\le x_n\le1/n\}$.
\end{ejercicio}
\begin{solucion}
El mapa toma valores en $H$ y es inyectivo. Para $y\in H$, la sucesión $x_n=y_n/n$ satisface $\sum_n|x_n|^2\le\sum_n1/n^2<\infty$ y pertenece a $K$, de modo que es sobreyectivo. Si $x^{(j)}\to x$ en $\ell^2$, entonces $|x_n^{(j)}-x_n|\le\norm{x^{(j)}-x}_2\to0$ para cada $n$; por tanto $nx_n^{(j)}\to nx_n$ y la convergencia coordenada implica convergencia en $H$. Así $T$ es continua. Como $K$ es compacto y $H$ es métrico, la biyección continua $T$ es un homeomorfismo. Obsérvese que no estamos afirmando continuidad de este multiplicador en todo $\ell^2$ con valores en $\ell^2$: aquí el codominio tiene la topología producto.
\end{solucion}

\begin{ejercicio}[Lectura 16, ejercicio 5]\label{ej:lectura-16-5}
Probar que todo espacio métrico separable es homeomorfo a un subespacio del toro infinito $\mathbb T^\infty=(S^1)^{\N}$.
\end{ejercicio}
\begin{solucion}
Por el teorema \ref{ch16:encaje}, basta encajar $H$ en $\mathbb T^\infty$. El mapa $h:[0,1]\to S^1$, $h(t)=e^{i\pi t}$, es continuo e inyectivo; por compacidad es un homeomorfismo sobre el semicírculo superior. Aplicamos $h$ coordenada a coordenada:
\[
\Phi((x_n))=(e^{i\pi x_n})_{n\ge1}.
\]
Este mapa es continuo e inyectivo en las topologías producto. Como $H$ es compacto y el producto numerable de circunferencias es métrico, $\Phi$ es un encaje. La composición $\Phi\circ J$ da el encaje de $X$. Se utiliza un semicírculo, no $e^{2\pi it}$ en todo $[0,1]$, pues este último identificaría los dos extremos.
\end{solucion}

\section{Encajes e imágenes continuas de compactos}
Una compactificación conserva la topología y añade límites; una inyección continua puede hacer algo distinto. El cubo de Hilbert contiene modelos de todos los espacios métricos separables, mientras que Cantor se proyecta continuamente sobre todos los compactos métricos no vacíos. Encajar y proyectar son operaciones diferentes, y distinguirlas evita errores en las demostraciones y en los ejercicios.
